\chapter{Les commandes utiles}
\label{ann:commandes}

\begin{chapeau}
Les commandes se lancent dans un terminal, depuis le dossier \fichier{CT-CORA3D/code_source/}, avec le Python du
projet (ici noté \code{python}). Sur Mac, c'est par exemple
\fichier{~/Library/Application Support/CTCORA3D/venv/bin/python}.
\end{chapeau}

\section{Le calcul sans interface}

\noindent\begin{minipage}{\linewidth}
\begin{lstlisting}[language=bash]
python lancer_tout.py --liste                  # inventaire des séries DICOM
python lancer_tout.py                          # toutes les séries, étapes 1 à 6
python lancer_tout.py --serie hip_to_knee      # une série (un morceau de nom unique suffit)
python lancer_tout.py --dossier "D:/CT chien 12" --sortie "D:/resultats" --seuil 250
python etape6_reperes.py ../exemples/greyhound_hindlimb_knee_to_ankle --os tibia
python etape7_guide.py --prescription prescription.json --os os.bin --sortie dossier \
                       --voisin fibula=fibula.bin --cartilage cartilage.u8
\end{lstlisting}
\end{minipage}

\section{Les tests et les campagnes}

\noindent\begin{minipage}{\linewidth}
\begin{lstlisting}[language=bash]
python tests/test_chaine.py                    # faux scanner (environ 3 s)
python tests/test_reperes.py                   # repères (environ 13 s)
python tests/test_guide.py                     # guides (environ 25 s ; --rapide)
cd cora && node tests/test_cas.mjs && node tests/test_reel.mjs && cd ..
node tests/campagne_strategie.mjs ../campagne_strategie 300 --exporter
python tests/campagne_guides.py ../campagne_strategie/guides --processus 4
node tests/campagne_export.mjs ../campagne tout 160 11
python tests/campagne_guides.py ../campagne --processus 4
\end{lstlisting}
\end{minipage}

Pour une campagne longue, limiter le nombre de fils par processus évite de surcharger le processeur :
\code{OMP_NUM_THREADS=3 ITK_GLOBAL_DEFAULT_NUMBER_OF_THREADS=3} avec \code{--processus 3}.

\section{Reconstruire}

\noindent\begin{minipage}{\linewidth}
\begin{lstlisting}[language=bash]
cd cora && python3 build_app.py && cd ..       # interface CORA après une modification du JavaScript
bash ../application/lanceurs/creer_app_autonome.sh       # application autonome macOS (environ 3 min)
\end{lstlisting}
\end{minipage}
Sous Windows, double-cliquer sur \texttt{application\textbackslash{}lanceurs\textbackslash{}creer\_app\_windows.bat} construit l'application
autonome (Python 3.10 ou plus récent et Internet nécessaires).

\section{Refaire ce document}

\noindent\begin{minipage}{\linewidth}
\begin{lstlisting}[language=bash]
cd documents/dossier_complet
bash construire.sh            # le PDF (trois passes de XeLaTeX)
bash construire.sh figures    # refait d'abord les figures et les schémas draw.io
\end{lstlisting}
\end{minipage}
